\documentclass[tikz,border=4pt]{standalone}
\usepackage{pe-tikz}
\begin{document}
\begin{tikzpicture}[x=0.42cm, y=0.62cm,
  tr/.style={fill=cA!55, draw=white, line width=0.6pt},
  te/.style={fill=cB!80, draw=white, line width=0.6pt},
  no/.style={fill=grid, draw=white, line width=0.6pt}]
  \foreach \k in {0,...,4}{
    \pgfmathtruncatemacro{\o}{12+3*\k}
    \foreach \i in {0,...,29}{
      \pgfmathtruncatemacro{\a}{\i<\o ? 1 : (\i<\o+3 ? 2 : 0)}
      \ifnum\a=1 \fill[tr] (\i,-\k) rectangle ++(1,0.8);\fi
      \ifnum\a=2 \fill[te] (\i,-\k) rectangle ++(1,0.8);\fi
      \ifnum\a=0 \fill[no] (\i,-\k) rectangle ++(1,0.8);\fi}
    \node[font=\scriptsize, anchor=east, text=ink] at (-0.3,-\k+0.4) {\pgfmathparse{int(\k+1)}\pgfmathresult. kezdőpont};}
  \draw[-{Stealth[length=2.2mm]}, line width=0.8pt, draw=ink] (0,-4.7) -- (30.6,-4.7) node[right, font=\small]{idő};
  \fill[tr] (2,1.4) rectangle ++(1,0.6); \node[font=\scriptsize, anchor=west] at (3.2,1.7) {tanítóadatok};
  \fill[te] (10,1.4) rectangle ++(1,0.6); \node[font=\scriptsize, anchor=west] at (11.2,1.7) {előrejelzési horizont ($h=3$)};
  \fill[no] (22,1.4) rectangle ++(1,0.6); \node[font=\scriptsize, anchor=west] at (23.2,1.7) {nem használt};
\end{tikzpicture}
\end{document}
